
% =============================================================================
% PROOFDOT 2 CHE DO
% =============================================================================
% Dung:
%   \begin{proofdot}{6}{a}\end{proofdot}  -> dong cham mot cot nhu ban cu
%   \begin{proofdot}{6}{b}\end{proofdot}  -> dong cham chia doi, co vach doc giua
%VD1: đang cần các câu từ 3 đến 9, chỉ cần gõ \printEX{{3}{9}}.
%+ VD2: đang cần các câu từ 3 đến 9, rồi câu 25 đến 37, chỉ cần gõ \printEX{{3}{9}{25}{37}}.
%+ VD3: đang cần các câu từ 3 đến 9, câu 18, rồi câu 21 đến 32, chỉ cần gõ \printEX{{3}{9}{18}{18}{21}{32}}.

\newcommand{\TSATheoremProofDotA}[1]{%
  \multido{}{#1}{\noindent\makebox[\linewidth]{\dotfill}\\[6.5mm]}%
}

% Cập nhật chế độ A (dòng chấm ngang) để sử dụng TikZ với độ đậm lớn hơn
\renewcommand{\TSATheoremProofDotA}[1]{%
	\ifnum#1>0\relax
	\par\addvspace{2pt}%
	\noindent\begin{tikzpicture}[x=\linewidth,y=1cm]
		% Giãn dòng bằng chế độ B (hệ số 1.1)
		\pgfmathsetmacro{\TSAProofDotLast}{#1*1.1}
		\pgfmathsetmacro{\TSAProofDotBottom}{-\TSAProofDotLast-.18}
		\useasboundingbox (0,.12) rectangle (1,\TSAProofDotBottom);
		% Kẻ dòng chấm ngang ĐẬM
		\foreach \i in {1,...,#1}{%
			\pgfmathsetmacro{\yy}{-\i*1.1}
			% Màu đen đậm và độ dày lớn hơn (.7pt)
			\draw[black!85,line width=.7pt,densely dotted] (.02,\yy) -- (.98,\yy);
		}%
	\end{tikzpicture}%
	\par\addvspace{2pt}%
	\fi
}

% Cập nhật chế độ B (chia cột) với độ đậm lớn hơn và giãn dòng 1.1
\providecommand{\TSATheoremProofDotB}{}
\renewcommand{\TSATheoremProofDotB}[1]{%
	\ifnum#1>0\relax
	\par\addvspace{1pt}%
	\noindent\begin{tikzpicture}[x=\linewidth,y=1cm]
		% Hệ số giãn dòng 1.1
		\pgfmathsetmacro{\TSAProofDotLast}{#1*1.1}
		\pgfmathsetmacro{\TSAProofDotBottom}{-\TSAProofDotLast-.18}
		\useasboundingbox (0,.12) rectangle (1,\TSAProofDotBottom);
		% Kẻ vạch dọc ở giữa (vẫn giữ xám để làm hướng dẫn nhẹ)
		\draw[gray!60,line width=.85pt] (.5,-.16) -- (.5,-\TSAProofDotLast-.12);
		\foreach \i in {1,...,#1}{%
			\pgfmathsetmacro{\yy}{-\i*1.1}
			% Màu đen đậm và độ dày lớn hơn (.7pt) cho các dòng chấm
			\draw[black!85,line width=.7pt,densely dotted] (.02,\yy) -- (.47,\yy);
			\draw[black!85,line width=.7pt,densely dotted] (.53,\yy) -- (.98,\yy);
		}%
	\end{tikzpicture}%
	\par\addvspace{2pt}%
	\fi
}

\newif\ifTSATheoremProofDotSplit
\newcommand{\TSATheoremProofDotSelect}[2]{%
  \TSATheoremProofDotSplitfalse
  \def\TSAProofDotMode{#2}%
  \def\TSAProofDotModeb{b}%
  \def\TSAProofDotModeB{B}%
  \ifx\TSAProofDotMode\TSAProofDotModeb\TSATheoremProofDotSplittrue\fi
  \ifx\TSAProofDotMode\TSAProofDotModeB\TSATheoremProofDotSplittrue\fi
  \ifTSATheoremProofDotSplit
    \TSATheoremProofDotB{#1}%
  \else
    \TSATheoremProofDotA{#1}%
  \fi
}

\makeatletter
% =============================================================================
% KIỂM TRA CHẾ ĐỘ ẨN / HIỆN CỦA PROOFDOT
% Tự động ẩn khi bật 1 trong 3 tùy chọn: loigiai, color, solcolor của ex_test
% =============================================================================
\newif\ifTSAHideProofDot
\def\TSAProofDotOverride{-1}% -1: tự động, 0: luôn hiện, 1: luôn ẩn
\newcommand{\HideProofDot}{\def\TSAProofDotOverride{1}}
\newcommand{\ShowProofDot}{\def\TSAProofDotOverride{0}}
\newcommand{\AutoProofDot}{\def\TSAProofDotOverride{-1}}

\newcommand{\TSACheckProofDotHide}{%
  \ifnum\TSAProofDotOverride=1\relax
    \TSAHideProofDottrue
  \else\ifnum\TSAProofDotOverride=0\relax
    \TSAHideProofDotfalse
  \else
    \TSAHideProofDotfalse
    % 1. Kiểm tra cờ \if@extest@hasanykey của ex_test (bật khi dùng loigiai, color, solcolor)
    \ifdefined\if@extest@hasanykey
      \if@extest@hasanykey \TSAHideProofDottrue \fi
    \fi
    % 2. Kiểm tra chuỗi tùy chọn \kindOPTN (loigiai, color, solcolor)
    \ifdefined\kindOPTN
      \IfSubStrEX{,solcolor,loigiai,color,}{,\kindOPTN,}{\TSAHideProofDottrue}{}%
    \fi
    % 3. Kiểm tra các cờ hiển thị lời giải / màu sắc
    \ifdefined\if@extest@showloigiai
      \if@extest@showloigiai \TSAHideProofDottrue \fi
    \fi
    \ifdefined\if@extest@showcolor
      \if@extest@showcolor \TSAHideProofDottrue \fi
    \fi
  \fi\fi
}

\DeclareDocumentEnvironment{proofdot}{m g +b}{%
  \TSACheckProofDotHide
  \ifTSAHideProofDot
    % Tự giác %: Ẩn hoàn toàn, không in dòng kẻ chấm hay sinh ra khoảng cách thừa
  \else
    % Chế độ in đề (dethi, book,...): In các dòng kẻ chấm
    \par\addvspace{-4pt}%
    \IfNoValueTF{#2}{\TSATheoremProofDotSelect{#1}{a}}{\TSATheoremProofDotSelect{#1}{#2}}%
    \par\addvspace{.5\topsep}%
  \fi
}{}
\makeatother
%%%%%
%Xài ký hiệu không chia hết, chia hết
\def\kch{\mathrel{\ooalign{$\vdots$\cr\hidewidth$\raisebox{0.25ex}{$/$}$\hidewidth\cr}}}
\def\ch{\mathrel{\vdots}}